\section{Data Type ও Field Properties}

\subsubsection{দাম ও ফোন নম্বর: দুটিতেই অঙ্ক, কাজ কি একই?}
দোকানের পণ্যের দাম রাখার Price ফিল্ডে ``দুইশ টাকা'' লেখা থাকলে সেই মান দিয়ে সরাসরি মোট মূল্য হিসাব করা কঠিন হয়।
অন্যদিকে যোগাযোগের নম্বরকে সাধারণ সংখ্যা হিসেবে রাখলে তার শুরুতে থাকা শূন্য হারিয়ে যেতে পারে। দাম এমন
সংখ্যামান হিসেবে রাখা প্রয়োজন, যার ওপর হিসাব করা যায়। যোগাযোগের নম্বর পরিচিতিমূলক লেখা হিসেবে রাখতে হবে,
যেন তার প্রতিটি অঙ্ক অক্ষুণ্ণ থাকে।

ব্যাংকের টাকার পরিমাণ, হাসপাতালের জন্মতারিখ ও ওয়েবসাইটের যোগাযোগ নম্বর ভিন্ন কাজে ব্যবহৃত হয়। টাইপ নির্বাচনে তাই মানের ব্যবহার বিবেচনা করতে হয়। শুধু টাইপ ঠিক রাখলেই তথ্যের সত্যতা যাচাইয়ের সব কাজ সম্পন্ন হয় না।

\begin{definitionbox}{Data Type}
একটি ফিল্ডে কী ধরনের মান রাখা যাবে এবং সেই মানের ওপর কী ধরনের কাজ করা যাবে তা নির্ধারণকারী বৈশিষ্ট্য হলো \textbf{Data Type}।
\end{definitionbox}


\subsubsection{Data Type নির্বাচনের গুরুত্ব}
মানের আকার ও প্রদর্শনের বিন্যাসও বিবেচ্য। Date/Time মান দিয়ে তারিখভিত্তিক কাজ করা যায়। একই তারিখ Text হিসেবে রাখলে তুলনা বা হিসাবের আগে উপযুক্ত রূপান্তরের প্রয়োজন হতে পারে।

\textbf{মান, ধরন ও প্রদর্শন আলাদা:} কোনো Number ফিল্ডে Currency-র মতো Format দিলেই সেটির Data Type Currency হয়ে যায় না। আবার ঠিক Data Type বাছলেও সব বাস্তব শর্ত পূরণ হয় না। GPA সংখ্যামূলক হওয়া এক শর্ত; নির্ধারিত সীমার মধ্যে থাকা আরেক শর্ত। তাই Type-এর সঙ্গে প্রয়োজনীয় Validation Rule-ও বসাতে হয়। নিচের ফিল্ড-প্রপার্টি আলোচনায় প্রতিটি নিয়ন্ত্রণের আলাদা কাজ দেখানো হয়েছে।


\subsection{Text ও Memo}
\textbf{Text বা Short Text} অক্ষর, সংখ্যা ও চিহ্নের সমন্বয়ে গঠিত ছোট মান রাখে। Name, Address, BloodGroup, Mobile ও PostalCode এর উদাহরণ। Access-এর Short Text-এ সর্বোচ্চ 255 ক্যারেক্টার রাখা যায়; Field Size দিয়ে প্রয়োজন অনুযায়ী আরও ছোট সীমা নির্ধারণ করা যায়। যেমন Name-এর সীমা 50 হলে 50 ক্যারেক্টারের বেশি নাম সরাসরি গ্রহণযোগ্য হবে না।

রোল 0012 একটি পরিচিতি হলে শুরুর শূন্য সংরক্ষণের জন্য Text ব্যবহার করা যায়। রোল দিয়ে গাণিতিক হিসাব দরকার কি না, তা দেখে ডেটা টাইপ নির্ধারণ করতে হয়।

\textbf{Memo বা Long Text} দীর্ঘ বিবরণ, মন্তব্য বা নির্দেশনা রাখার জন্য ব্যবহৃত হয়। পুরোনো বই ও Access-এ Text এবং Memo নাম দেখা যায়; নতুন সংস্করণে Short Text ও Long Text নাম প্রচলিত। Memo-এর জন্য পুরোনো সংস্করণভিত্তিক প্রশ্নে প্রায় 65,535 ক্যারেক্টারের সীমা উল্লেখ থাকে; আধুনিক Long Text-এ তথ্য সংরক্ষণের সীমা ও কন্ট্রোলে তথ্য প্রদর্শনের সীমা এক নয়। তাই সংস্করণ উল্লেখ থাকলে তার নির্দিষ্ট নিয়ম অনুসরণ করতে হবে।

\textbf{উদাহরণ:} Roll যদি 001 হয়, Text হিসেবে তার তিনটি চিহ্ন বজায় থাকে। Number হিসেবে একই মান 1 হয়। GPA 4.75-এর ওপর তুলনা ও গড় নির্ণয় প্রয়োজন, তাই তা সংখ্যাভিত্তিক টাইপে রাখা যুক্তিযুক্ত। পরিচিতির মান এবং গণনার রাশি আলাদা করে বুঝতে হবে।

\subsection{Number ও তার Field Size}
পূর্ণসংখ্যা বা ভগ্নাংশযুক্ত সংখ্যামান রাখার জন্য \textbf{Number} ব্যবহৃত হয়। Age, Marks, Quantity ও GPA এর উদাহরণ। Access-এ Number-এর Field Size শুধু প্রদর্শনের দৈর্ঘ্য নয়; এটি গ্রহণযোগ্য সংখ্যার সীমা, ভগ্নাংশের সুবিধা এবং সংরক্ষণের আকার নির্ধারণ করে।
\begin{flushleft}\small
\begin{tabularx}{\linewidth}{|p{28mm}|p{16mm}|X|}\hline
\chSixHeader Field Size & আকার & গ্রহণযোগ্য মান/বৈশিষ্ট্য\\\hline
Byte & 1 বাইট & 0 থেকে 255; পূর্ণসংখ্যা\\\hline
Integer & 2 বাইট & $-32768$ থেকে $32767$; পূর্ণসংখ্যা\\\hline
Long Integer & 4 বাইট & $-2147483648$ থেকে $2147483647$; পূর্ণসংখ্যা\\\hline
Single & 4 বাইট & ভাসমান দশমিক; প্রায় 7 গুরুত্বপূর্ণ অঙ্ক\\\hline
Double & 8 বাইট & ভাসমান দশমিক; প্রায় 15 গুরুত্বপূর্ণ অঙ্ক\\\hline
Decimal & 12 বাইট & নির্ধারিত precision ও scale-সহ দশমিক\\\hline
\end{tabularx}
\end{flushleft}
Age-এর ছোট পূর্ণসংখ্যার জন্য Byte যথেষ্ট হতে পারে, কিন্তু GPA-এর ভগ্নাংশ রাখতে Integer উপযুক্ত নয়। বড় গণনায় Double বেছে নেওয়া যায়। প্রদর্শনে Decimal Places পরিবর্তন করা আর পূর্ণসংখ্যার Field Size-কে ভগ্নাংশ গ্রহণে সক্ষম করা এক কাজ নয়। ভাসমান দশমিকের হিসাবে আনুমানিকতার বিষয়ও থাকে।

\subsection{AutoNumber, Currency ও Yes/No}
\textbf{AutoNumber} নতুন রেকর্ডের জন্য স্বয়ংক্রিয় পরিচিতি তৈরি করে। এটি প্রাইমারি কী তৈরিতে সুবিধাজনক। রেকর্ড বাদ পড়লে বা কোনো এন্ট্রি বাতিল হলে ধারায় ফাঁক থাকতে পারে; তাই AutoNumber-কে মোট রেকর্ডের সংখ্যা হিসেবে ধরা যায় না। প্রচলিত Long Integer AutoNumber-এর সঙ্গে সম্পর্ক তৈরির জন্য অন্য টেবিলের ফরেন কী সাধারণত Number--Long Integer রাখা হয়। সেখানে আবার AutoNumber দিলে একই ব্যক্তির সঠিক পরিচিতি স্বয়ংক্রিয়ভাবে মিলবে না।

টাকার পরিমাণ সংরক্ষণে \textbf{Currency} ব্যবহৃত হয়। Access-এ এটি 8 বাইটের স্থির দশমিক মান এবং দশমিকের পরে চার ঘর পর্যন্ত নির্ভুলতা সমর্থন করে। বেতন, ফি ও পণ্যের দাম এর উদাহরণ। মুদ্রার চিহ্ন বা কমা দেখানো Format-এর কাজ; সংরক্ষিত মানের অর্থ আলাদা।

\textbf{Yes/No} দুটি অবস্থার একটি প্রকাশ করে; যেমন Paid, Active বা Present। পর্দায় Yes/No, True/False, On/Off বা Check Box দেখা যেতে পারে। Access-এ True-এর অভ্যন্তরীণ সংখ্যামান $-1$ এবং False-এর 0; সব সফটওয়্যারে একই অভ্যন্তরীণ রূপ ধরে নেওয়া যাবে না।

\subsection{Date/Time ও তার বিন্যাস}
জন্মতারিখ, যোগদানের তারিখ, সাক্ষাতের সময় বা লেনদেনের সময় রাখতে \textbf{Date/Time} ব্যবহৃত হয়। এতে শুধু লেখা হিসেবে তারিখ রাখার তুলনায় সময়ের ক্রমে সাজানো ও তারিখভিত্তিক হিসাব করা সহজ হয়। প্রচলিত Access Date/Time 8 বাইট ব্যবহার করে।
\begin{flushleft}\small
\begin{tabularx}{\linewidth}{|p{35mm}|X|}\hline
\chSixHeader Format & প্রদর্শনের ধরন\\\hline
General Date & প্রয়োজন অনুযায়ী তারিখ ও সময়\\\hline
Short Date & সংক্ষিপ্ত সংখ্যায় তারিখ, যেমন 15/09/2026\\\hline
Medium Date & মাসের সংক্ষিপ্ত নামসহ তারিখ\\\hline
Long Date & বার/মাসের পূর্ণ নামসহ বিস্তৃত তারিখ\\\hline
Short/Medium/Long Time & ঘণ্টা-মিনিট, প্রয়োজনে সেকেন্ড বা AM/PM\\\hline
\end{tabularx}
\end{flushleft}
অঞ্চল ও সংস্করণভেদে প্রদর্শনের রূপ বদলাতে পারে। একই তারিখ 15/09/2026 বা 15-Sep-2026 দেখালেও Format বদলানোয় মূল তারিখের অর্থ বদলায় না। 03/04/2026-এর মতো দ্ব্যর্থক তারিখ এন্ট্রির সময় দিনের ও মাসের নিয়ম নিশ্চিত করা প্রয়োজন।

\subsection{OLE Object, Attachment, Hyperlink ও Lookup}
OLE-এর পূর্ণরূপ \textbf{Object Linking and Embedding}। \textbf{OLE Object} দিয়ে অন্য অ্যাপ্লিকেশনের অবজেক্ট, যেমন ছবি বা নথি সংযুক্ত/এম্বেড করার পুরোনো Access পদ্ধতি বোঝায়। আধুনিক \textbf{Attachment} টাইপেও একটি রেকর্ডের সঙ্গে ছবি ও নথি রাখা যায়; এটি পুরোনো \texttt{.mdb} ফরম্যাটে নেই। প্রশ্নে পুরোনো Access-এর Photo ফিল্ড থাকলে OLE Object প্রত্যাশিত হতে পারে, আর আধুনিক নকশায় Attachment প্রাসঙ্গিক।

\textbf{Hyperlink} দিয়ে ওয়েব ঠিকানা, ফাইল বা অন্যান্য গন্তব্যের লিংক রাখা হয়। এটি ছবির বাইনারি ডেটা সংরক্ষণের সমার্থক নয়। \textbf{Lookup Wizard} দিয়ে টেবিল, কুয়েরি বা নির্দিষ্ট তালিকা থেকে মান বাছাইয়ের ব্যবস্থা তৈরি করা যায়। যেমন বিভাগ হাতে না লিখে তালিকা থেকে নির্বাচন করা যায়। Lookup নিজে পৃথক মৌলিক সংরক্ষণ টাইপ নয়; অন্তর্নিহিত ফিল্ড Text বা Number হতে পারে।

\chSixInterfaceSide{ch06_access_data_types.png}{Design View-তে Data Type নির্বাচন}{Data Type-এর তালিকা থেকে ফিল্ডে কী ধরনের মান রাখা হবে তা নির্বাচন করা হয়। যেমন GPA-এর জন্য Number এবং জন্মতারিখের জন্য Date/Time উপযোগী। নিচের বৈশিষ্ট্যগুলো নির্বাচিত টাইপ অনুসারে বদলাতে পারে।}

\subsection{Field Properties}
ডেটা টাইপের সঙ্গে ফিল্ডের অন্যান্য বৈশিষ্ট্য নির্ধারণ করলে এন্ট্রি ও উপস্থাপন নিয়ন্ত্রিত হয়। সব টাইপে সব বৈশিষ্ট্য প্রযোজ্য নয়। নিচের উদাহরণগুলো Access-এর ফিল্ড নকশা বোঝায়।
\begin{flushleft}\small
\begin{tabularx}{\linewidth}{|p{31mm}|X|}\hline
\chSixHeader Property & কাজ ও উদাহরণ\\\hline
Field Size & Text-এর সর্বোচ্চ দৈর্ঘ্য বা Number-এর সীমা/নির্ভুলতা\\\hline
Format & মানের প্রদর্শনরীতি; যেমন মুদ্রা বা তারিখের বিন্যাস\\\hline
Input Mask & এন্ট্রির ছক; যেমন নির্দিষ্ট স্থানে অঙ্ক লেখা\\\hline
Caption & দৃশ্যমান শিরোনাম; Name-এর বদলে শিক্ষার্থীর নাম দেখানো\\\hline
Default Value & নতুন রেকর্ডে প্রাথমিক মান; যেমন City-তে Dhaka\\\hline
Validation Rule & গ্রহণের শর্ত; GPA-এর জন্য \texttt{Between 0 And 5}\\\hline
Validation Text & নিয়ম ভাঙলে অর্থবহ বার্তা দেখানো\\\hline
Required & মান অবশ্যই দিতে হবে কি না; Yes হলে NULL নয়\\\hline
Allow Zero Length & Text-এ শূন্য দৈর্ঘ্যের স্ট্রিং গ্রহণযোগ্য কি না\\\hline
Indexed & ইনডেক্স হবে কি না এবং একই মান পুনরাবৃত্ত হতে পারবে কি না\\\hline
\end{tabularx}
\end{flushleft}
\textbf{Format ও Input Mask:} Format সাধারণত মান কীভাবে দেখাবে তা নিয়ন্ত্রণ করে, Input Mask কী ছকে এন্ট্রি হবে তা নিয়ন্ত্রণ করে। মাসিক ফি 1500-কে মুদ্রা চিহ্নসহ দেখানো Format-এর কাজ। মোবাইল ফিল্ডে নির্দিষ্ট সংখ্যক অঙ্ক লিখতে সাহায্য করা Mask-এর কাজ। Mask মেনে লেখা নম্বর বাস্তবে সঠিক ব্যক্তির কি না, তা অবশ্য পৃথক যাচাইয়ের বিষয়।

\textbf{Default ও Validation:} Default Value শুধু প্রাথমিক প্রস্তাবিত মান; অনুমতি থাকলে ব্যবহারকারী তা বদলাতে পারে। Validation Rule গ্রহণযোগ্যতার শর্ত। GPA-এর Default 0 রাখা মানে সব রেকর্ডের GPA সর্বদা 0 থাকবে না; আর \texttt{Between 0 And 5} দিলে 6 গ্রহণ করা যাবে না। Validation Text-এ ``GPA 0 থেকে 5-এর মধ্যে লিখুন'' দিলে ব্যবহারকারী ভুলের কারণ বুঝতে পারে।


\textbf{টাইপ, ডোমেইন ও সত্যতা:} GPA-এর জন্য Number হলো মানের মৌলিক ধরন; 0 থেকে 5-এর অনুমোদিত সীমা তার Domain-এর নিয়ম। 6.20 সংখ্যাই, তবু এই নিয়মে অবৈধ। আবার আসল GPA 4.20-এর বদলে 4.80 লেখা হলে টাইপ ও সীমা উভয় পরীক্ষায় পাস করবে, কিন্তু তথ্য সত্য হবে না। Validation ভুলের একটি অংশ আটকায়; মূল ফলপত্র মিলিয়ে দেখা জরুরি।

\subsubsection{Input Mask, Caption ও Default Value-এর প্রয়োগ}
Access-এর Input Mask-এ \texttt{0} সাধারণত আবশ্যিক অঙ্ক এবং \texttt{9} ঐচ্ছিক অঙ্ক বা ফাঁকা স্থানের জায়গা নির্দেশ করে। \texttt{L} আবশ্যিক অক্ষর, \texttt{?} ঐচ্ছিক অক্ষর, \texttt{A} আবশ্যিক অক্ষর/অঙ্ক এবং \texttt{a} ঐচ্ছিক অক্ষর/অঙ্কের ধারণা প্রকাশ করে। মাস্ক \texttt{00000} পাঁচটি অঙ্কের ছক দিতে পারে। এতে 12 লিখে এন্ট্রি সম্পূর্ণ করা যাবে না, কিন্তু পাঁচটি ভুল অঙ্কও ছক মেনে যেতে পারে।

\textbf{Caption-এর উদাহরণ:} টেবিলে ফিল্ডের নাম StudentName রেখে পর্দায় ``শিক্ষার্থীর নাম'' দেখানো যায়। SQL লেখার সময় সাধারণত প্রকৃত ফিল্ডের নাম ব্যবহার করতে হয়। দৃশ্যমান Caption দেখে সেটিকেই টেবিলের নতুন ফিল্ডের নাম ধরে নেওয়া ভুল।

\textbf{Default Value-এর উদাহরণ:} নতুন রেকর্ডে GroupName-এর প্রাথমিক মান Science রাখা হলো। Humanities-এর শিক্ষার্থী হলে সেই মান বদলাতে হবে। Default Value সুবিধা দেয়, কিন্তু বাস্তব তথ্য যাচাই ছাড়া অন্ধভাবে গ্রহণ করলে ভুল তৈরি হয়।

\subsubsection{Required ও Allow Zero Length-এর সমন্বয়}
Text ফিল্ডে NULL এবং শূন্য দৈর্ঘ্যের স্ট্রিং আলাদা। নিচের সারণি শুধু Required ও Allow Zero Length-এর সমন্বয়ের ধারণা দেখায়; বাড়তি Validation Rule থাকলে আরও সীমা প্রযোজ্য হতে পারে।
\begin{flushleft}\small
\begin{tabular}{|l|l|l|l|}\hline
\chSixHeader Required & Allow Zero Length & NULL & শূন্য দৈর্ঘ্য\\\hline
No & No & অনুমোদিত & অননুমোদিত\\\hline
No & Yes & অনুমোদিত & অনুমোদিত\\\hline
Yes & No & অননুমোদিত & অননুমোদিত\\\hline
Yes & Yes & অননুমোদিত & অনুমোদিত\\\hline
\end{tabular}
\end{flushleft}
শুধু কয়েকটি ফাঁকা স্থান লেখা শূন্য দৈর্ঘ্য নয়। অর্থবহ নাম বাধ্যতামূলক করতে প্রয়োজনে ফাঁকা স্থান বাদ দিয়ে দৈর্ঘ্য যাচাই করতে হয়। Validation Rule ডেটার গ্রহণযোগ্যতা এবং Validation Text সেই নিয়মভঙ্গের বার্তা প্রকাশ করে; দুটি আলাদা বৈশিষ্ট্য।


\subsection{NULL, শূন্য ও খালি স্ট্রিং}
অজানা বা অনুপস্থিত মান বোঝাতে \textbf{NULL} ব্যবহৃত হয়। পরীক্ষায় নম্বর 0 পাওয়া এবং নম্বর জানা না থাকা এক নয়। একইভাবে Text-এর \texttt{''} একটি শূন্য দৈর্ঘ্যের স্ট্রিং; NULL কোনো সাধারণ স্ট্রিং নয়। SQL-এ NULL খুঁজতে \texttt{IS NULL}, আর অ-NULL মান খুঁজতে \texttt{IS NOT NULL} লেখা হয়। \texttt{= NULL} দিয়ে স্বাভাবিক মানের মতো তুলনা করা যায় না।

Required=Yes বা NOT NULL অনুপস্থিত মান প্রতিরোধ করে; তা নিজে সব ক্ষেত্রে খালি স্ট্রিং বা শুধু ফাঁকা স্থান প্রতিরোধ করে না। অর্থবহ নাম বাধ্যতামূলক করতে প্রয়োজন হলে শূন্য দৈর্ঘ্য ও ফাঁকা স্থানের জন্য অতিরিক্ত যাচাই লাগে।

\subsection{SQL-এ লেখা ও সংখ্যার টাইপের নাম}
Access-এর পর্দায় ব্যবহৃত টাইপের নাম এবং SQL নির্দেশে ঘোষিত নাম সবসময় এক নয়। \texttt{INT} বা \texttt{INTEGER} পূর্ণসংখ্যার জন্য ব্যবহৃত হয়; \texttt{SMALLINT} কম সীমার পূর্ণসংখ্যার টাইপ। \texttt{REAL} ও \texttt{DOUBLE PRECISION} ভাসমান দশমিকের নাম। টাইপের সমর্থন, আকার ও সীমা DBMS অনুযায়ী জানতে হয়।

\texttt{CHAR(n)} নির্দিষ্ট দৈর্ঘ্যের এবং \texttt{VARCHAR(n)} সর্বোচ্চ নির্ধারিত দৈর্ঘ্যের পরিবর্তনশীল লেখা রাখার জন্য ব্যবহৃত হয়। যেমন PostgreSQL-এ \texttt{CHAR(5)}-এ ছোট লেখা রাখলে অবশিষ্ট জায়গা স্পেস দিয়ে পূরণ হয়; \texttt{VARCHAR(5)}-এ ছোট লেখার জন্য এমন পূরণ প্রয়োজন হয় না। এখানে 5 অক্ষরের সীমা, সংরক্ষণের মোট বাইট নয়। SQLite-এ \texttt{VARCHAR(5)} ঘোষণা করলেই পাঁচ অক্ষরের সীমা কার্যকর হয় না।

নির্দিষ্ট দশমিক নির্ভুলতা দরকার হলে উপযুক্ত DBMS-এ \texttt{DECIMAL(p,s)} বা \texttt{NUMERIC(p,s)} ব্যবহার করা যায়। এখানে $p$ মোট অঙ্কের সংখ্যা এবং $s$ দশমিকের পরের অঙ্কের সংখ্যা। PostgreSQL-এ \texttt{NUMERIC(8,2)} মোট আট অঙ্কের মধ্যে দশমিকের পরে দুই অঙ্কের সীমা নির্ধারণ করে। এটি ভাসমান দশমিকের টাইপ নয়। SQLite-এর NUMERIC affinity-কে এই ধরনের নির্দিষ্ট দশমিক বিধির সমার্থক ধরা যাবে না।

\subsection{জেনে রাখো: Access ও SQLite-এর টাইপের পার্থক্য}
SQLite-এর সাধারণ টেবিলে মানের সংরক্ষণশ্রেণি হলো NULL, INTEGER, REAL, TEXT ও \textbf{BLOB (Binary Large Object)}। INTEGER পূর্ণসংখ্যা, REAL ভাসমান দশমিক, TEXT অক্ষরসমষ্টি এবং BLOB বাইনারি উপাত্ত রাখে। কলামের ঘোষিত টাইপ থেকে TEXT, NUMERIC, INTEGER, REAL বা BLOB affinity নির্ধারিত হতে পারে। NUMERIC affinity ও স্বতন্ত্র storage class একই বিষয় নয়।

SQLite-এ Access-এর মতো পৃথক Date/Time বা Currency storage class নেই। ধারাবাহিক \texttt{YYYY-MM-DD} বিন্যাসে TEXT দিয়ে তারিখ রাখা যায়। Boolean-এ প্রচলিতভাবে 0 ও 1 ব্যবহৃত হয়। নির্দিষ্ট টাইপ ও মান নিশ্চিত করতে CHECK, NOT NULL বা উপযুক্ত নকশা প্রয়োজন; সাধারণ SQLite টেবিলের টাইপ-আচরণকে Access-এর সঙ্গে পুরোপুরি এক ধরা যাবে না।

\subsection{সংরক্ষণের আকার ও সংস্করণভিত্তিক তথ্য}
Number-এর Single-এ প্রায় $\pm3.4\times10^{38}$ এবং Double-এ প্রায় $\pm1.797\times10^{308}$ পর্যন্ত বড় মান প্রকাশ করা যায়। তবে সংখ্যার সীমা এবং নির্ভুলভাবে প্রকাশযোগ্য গুরুত্বপূর্ণ অঙ্কের সংখ্যা আলাদা বিষয়। বড় মান রাখা গেলেই প্রতিটি দশমিক নিখুঁতভাবে রাখা যাবে না। Replication ID একটি বিশেষ 16 বাইটের পরিচিতি, যা পুরোনো replicated database-এ ব্যবহৃত হয়। আধুনিক \texttt{.accdb}-তে পুরোনো replication ব্যবস্থা নেই।

সাধারণ AutoNumber-এর Long Integer রূপ 4 বাইট এবং বিশেষ Replication ID রূপ 16 বাইট হতে পারে। আগে পরিচিত OLE ও BLOB নাম দুটি বিভিন্ন সফটওয়্যারের সংরক্ষণ ও ব্যবহারের প্রসঙ্গে আসে। ছবি, শব্দ বা নথিও ডেটা হতে পারে---শুধু ফাইলের নাম ও প্রকৃত বাইনারি কনটেন্ট এক নয়।

পুরোনো পাঠ্যপুস্তকে Memo-এর জন্য 65,536 ক্যারেক্টার বা প্রায় 64K-এর হিসাব এবং কিছু সফটওয়্যার-সীমায় 65,535 দেখা যায়। এগুলোকে সব সংস্করণের সার্বজনীন Long Text সীমা হিসেবে ব্যবহার করা যাবে না। একইভাবে logical মানের ধারণাগত বিটসংখ্যা এবং কোনো সফটওয়্যারের ফিল্ডে ব্যবহৃত স্থান আলাদা হতে পারে। নির্দিষ্ট পুরোনো প্রশ্নে দেওয়া সফটওয়্যার ও সংস্করণের প্রেক্ষিত বিবেচনা করা প্রয়োজন।

\subsection{একটি কর্মী-সারণির জন্য Data Type নির্বাচন}
একটি অফিসে কর্মীর পরিচিতি, যোগাযোগ, বেতন, যোগদানের তারিখ, ছবি ও মন্তব্য রাখতে হবে। একই ব্যক্তি সম্পর্কে তথ্য হলেও প্রতিটি ফিল্ডের কাজ আলাদা। নিচের কাল্পনিক মানগুলো দেখে Access-এ উপযুক্ত টাইপ নির্বাচন করা যায়।
\begin{flushleft}\small
\begin{tabularx}{\linewidth}{|p{24mm}|p{31mm}|X|}\hline
\chSixHeader ফিল্ড ও নমুনা & টাইপ & নির্বাচনের কারণ\\\hline
EmpID: E021 & Short Text & পরিচিতিতে অক্ষর আছে; হিসাবের সংখ্যা নয়\\\hline
Name: Rima & Short Text & সীমিত দৈর্ঘ্যের নাম\\\hline
Mobile: 01700000321 & Short Text & শুরুতে শূন্যসহ পরিচিতি বজায় রাখা\\\hline
Salary: 18500.50 & Currency & টাকার পরিমাণে নির্দিষ্ট দশমিক প্রয়োজন\\\hline
JoinDate: 15-01-2025 & Date/Time & তারিখের অর্থে তুলনা ও সর্টিং\\\hline
Active: Yes & Yes/No & দুটি সম্ভাব্য অবস্থা\\\hline
Photo: কর্মীর ছবি & OLE Object/Attachment & সংশ্লিষ্ট অবজেক্ট বা ফাইল রাখা\\\hline
Remarks: দীর্ঘ বিবরণ & Long Text/Memo & সাধারণ নামের তুলনায় দীর্ঘ লেখা\\\hline
\end{tabularx}
\end{flushleft}
একই নকশা SQLite-এ করলে Text, Integer, Real, Numeric বা Blob-এর ভূমিকা অনুযায়ী ঘোষণা করতে হয়; Access-এর Currency বা OLE Object-এর নাম সরাসরি সেখানে ব্যবহার করার প্রয়োজন নেই। একটি সফটওয়্যারের ফিল্ডের আকার অন্যটির ক্ষেত্রেও একই হবে ধরে নেওয়া যাবে না।

ফিল্ডের ধারণক্ষমতা তার বর্তমান মানের দৈর্ঘ্য নয়। Name-এর Field Size 30 হলে সর্বোচ্চ 30 অক্ষরের নাম গ্রহণের সীমা নির্ধারিত হয়। বর্তমান নাম Rima চার অক্ষরের হলেও ভবিষ্যতে দীর্ঘ নাম আসতে পারে; তাই শুধু Rima দেখে সীমা 4 করা উচিত নয়। মোট ডেটাবেজ ফাইলের আকারও দৃশ্যমান অক্ষর গুনে পাওয়া যায় না, কারণ অন্যান্য কাঠামো সংরক্ষণেও জায়গা লাগে।

\subsection{গুরুত্বপূর্ণ প্রশ্ন}
\textbf{অনুশীলন:} Name, Salary, BirthDate, Photo, Active ও Mobile-এর Access টাইপ নির্বাচন করো। GPA-এর জন্য Integer কেন অনুপযুক্ত এবং Required=Yes থাকলেও খালি স্ট্রিং আলাদাভাবে যাচাই করতে হয় কেন, ব্যাখ্যা করো।


